\chapter{Evaluation and validation}
\label{ch:evaluation}

\section{Evaluation strategy}

The evaluation in this thesis is centered on ecological explanatory fit for the 2005--2021 window, using scenario-based counterfactual analysis.

The reference annual trajectory is the Romanian national series (2005--2021), from 5800 bears in 2005 to 7500 bears in 2021. This implies an observed cumulative increase of 29.31\% and an average annual growth of approximately 1.62\%/year.

The simulations are compared against this trajectory with:

\begin{itemize}
    \item MAE (mean absolute error),
    \item RMSE (root mean square error),
    \item end-of-window gap to 2021 reference,
    \item cumulative growth and annualized growth inside the same window.
\end{itemize}

\section{Experimental protocol and scenarios}

All scenarios were run with three replicates and then post-processed into annual aggregates (year 0 to year 16, mapped to 2005 to 2021). The campaign includes:

\begin{itemize}
    \item baseline,
    \item EU protection driver (2007),
    \item hunting-ban driver (2016),
    \item intervention-cull driver (2021),
    \item habitat-loss driver,
    \item anthropogenic-food driver,
    \item combined realistic schedule.
\end{itemize}

\begin{table}[htbp]
\centering
\caption{Scenario definitions used in the final campaign.}
\label{tab:ch4_scenario_matrix}
\begin{tabular}{|l|p{8.5cm}|}
\hline
Scenario & Main schedule effect \\
\hline
Baseline & No time-varying policy/habitat driver beyond baseline settings. \\
EU protection 2007 & Danger mortality reduced from simulation year 2. \\
Hunting ban 2016 & Danger mortality reduced from simulation year 11. \\
Intervention cull 2021 & Danger mortality increased from simulation year 16. \\
Habitat loss & Forest food reduced stepwise (years 3, 9, 14). \\
Anthropogenic food & Food regrowth increased stepwise (years 5, 10, 15). \\
Combined realistic & Layered schedule: protection + food increase + 2016 ban + 2021 cull. \\
\hline
\end{tabular}
\end{table}

\section{Ecological results}

\subsection{Quantitative summary}

Table~\ref{tab:ch4_ecology_summary} shows the core quantitative comparison against the 2005--2021 reference trajectory.

\begin{table}[htbp]
\centering
\caption{Ecological validation metrics (annual aggregates, 2005--2021).}
\label{tab:ch4_ecology_summary}
\begin{tabular}{|l|c|c|c|c|c|}
\hline
Scenario & End pop. & Growth (\%) & CAGR (\%/y) & MAE & RMSE \\
\hline
Baseline & 6020.67 & 2.57 & 0.16 & 360.32 & 539.66 \\
EU protection 2007 & 6103.33 & 3.98 & 0.24 & 342.60 & 511.44 \\
Hunting ban 2016 & 6068.67 & 3.39 & 0.21 & 348.76 & 522.58 \\
Intervention cull 2021 (16y) & 6020.67 & 2.57 & 0.16 & 360.32 & 539.66 \\
Habitat loss & 6020.67 & 2.57 & 0.16 & 360.32 & 539.66 \\
Anthropogenic food & 7175.00 & 22.24 & 1.26 & 261.25 & 336.27 \\
Combined realistic & 7194.67 & 22.58 & 1.28 & 278.64 & 344.65 \\
\hline
Reference (observed) & 7500.00 & 29.31 & 1.62 & 0.00 & 0.00 \\
\hline
\end{tabular}
\end{table}

Main observations:

\begin{enumerate}
    \item Baseline and policy-only schedules underpredict growth strongly in this calibration.
    \item The anthropogenic-food schedule yields the best fit in this campaign (lowest MAE and RMSE).
    \item The combined realistic schedule is very close to anthropogenic food, but slightly worse on MAE/RMSE.
    \item The 2021 intervention-cull schedule has negligible effect within the 2005--2021 comparison window because it is activated at the end of the window.
\end{enumerate}

Relative to baseline, MAE drops by about 27.5\% under the anthropogenic-food scenario and by about 22.7\% under the combined scenario.

\subsection{End-of-window interpretation}

At year 16 (mapped to 2021), the reference value is 7500 bears.

\begin{itemize}
    \item Baseline gap: 1479.33 bears below reference.
    \item Anthropogenic-food gap: 325.00 bears below reference.
    \item Combined realistic gap: 305.33 bears below reference.
\end{itemize}

This indicates that, in the current parameterization, increased food availability is the dominant mechanism needed to approach the observed 2021 level.

\subsection{Mortality-structure summary}

Across scenarios, mean annual total deaths decrease from approximately 512.76 (baseline) to about 491.47--491.61 in the two best-fit scenarios. The danger-death share also decreases in the combined schedule relative to baseline, consistent with lower danger mortality settings in part of the timeline.

\section{Figure placeholders for final insertion}

The following placeholders are intentionally left as boxes so final plots can be inserted after figure export.

\begin{figure}[htbp]
    \centering
    % Suggested final file: figures/ch4_population_all_scenarios.pdf
    \fbox{\rule{0pt}{2.2in}\rule{0.85\textwidth}{0pt}}
    \caption{Observed reference trajectory (2005--2021) versus all simulated scenario trajectories.}
    \label{fig:ch4_population_all}
\end{figure}

\begin{figure}[htbp]
    \centering
    % Suggested final file: figures/ch4_population_top3.pdf
    \fbox{\rule{0pt}{2.2in}\rule{0.85\textwidth}{0pt}}
    \caption{Focused comparison of baseline, anthropogenic-food, and combined-realistic scenarios against observed data.}
    \label{fig:ch4_population_top3}
\end{figure}

\begin{figure}[htbp]
    \centering
    % Suggested final file: figures/ch4_residuals_by_year.pdf
    \fbox{\rule{0pt}{2.2in}\rule{0.85\textwidth}{0pt}}
    \caption{Year-by-year residuals (simulated minus observed) for selected scenarios.}
    \label{fig:ch4_residuals}
\end{figure}

\begin{figure}[htbp]
    \centering
    % Suggested final file: figures/ch4_deaths_components_selected.pdf
    \fbox{\rule{0pt}{2.2in}\rule{0.85\textwidth}{0pt}}
    \caption{Deaths by cause (danger, starvation, old age) for baseline and best-fit scenarios.}
    \label{fig:ch4_death_components}
\end{figure}

\begin{figure}[htbp]
    \centering
    % Suggested final file: figures/ch4_births_vs_deaths_selected.pdf
    \fbox{\rule{0pt}{2.2in}\rule{0.85\textwidth}{0pt}}
    \caption{Births versus total deaths over time for baseline and best-fit scenarios.}
    \label{fig:ch4_births_deaths}
\end{figure}

\section{Computational validation note}

The current chapter uses the newly generated ecological campaign data from the Results directory. Detailed runtime benchmarking (sequential vs multithreaded vs distributed) is discussed in Chapter 3 and can be appended here with the same table/figure pattern if a dedicated benchmark rerun is performed on the final hardware.

\section{Interpretation and validity limits}

The most important conclusion from this campaign is comparative: in the tested schedules, food-access increase explains substantially more of the 2005--2021 rise than policy-only danger-mortality shifts.

Two limits are important:

\begin{enumerate}
    \item The folders marked as weird results indicate the need for an additional reproducibility check before publication-grade claims about driver ranking.
    \item The intervention-cull effect should be evaluated on a longer post-2021 projection horizon, because activation at year 16 leaves almost no in-window response time.
\end{enumerate}

Even with these limits, the campaign is strong enough to support a final thesis-level conclusion about which mechanisms are necessary to approach observed growth in the validated window.
